\chapter{Advanced Legacy Proxy System}
\label{ch:alps}

Die identifizierten Limitationen von \gls{lps} in Abschnitt~\ref{sec:lps-limitationen} motivieren die Entwicklung eines neuen Systems, das diese Limitationen adressiert und gleichzeitig die Vorteile von \gls{lps} beibehält. Das neue System wird als \gls{alps} realisiert und in diesem Kapitel vorgestellt.

\section{Anforderungen}
\label{sec:alps-anforderungen}

Die Anforderung an \gls{alps} lassen sich direkt aus den Vorteilen und Limitationen von \gls{lps} ableiten:

\begin{itemize}
    \item \textbf{Asynchrone Entkopplung beibehalten}: Das vorteilhafte asynchrone Design von \gls{lps} soll beibehalten werden, da es sich positiv auf die Antwortzeiten des Systems auswirkt, wie in Abschnitt~\ref{sub-sec:lps-auswirkungen} beschrieben wurde.
    \item \textbf{Transparenz zum Legacy-System beibehalten}: Es soll möglich sein, \gls{alps} in das bestehende Legacy-System transparent zu integrieren.
    \item \textbf{Nachrichtenverlust verbieten}: Alarmmeldungen dürfen nicht verloren gehen.
    \item \textbf{Exactly-Once-Delivery garantieren}: Alarmmeldungen dürfen nicht mehrfach verarbeitet werden.
    \item \textbf{In-Order-Delivery pro Alarmanlage garantieren}: Alarmmeldungen einzelner Alarmanlagen sollen in der richtigen Reihenfolge verarbeitet werden.
    \item \textbf{Vendor-Neutralität}: \gls{alps} soll nicht von einem bestimmten Vendor abhängig sein, um die Portabilität und Flexibilität von \gls{alps} zu ge\-währ\-leis\-ten.
\end{itemize}

\section{Architektur}
\label{sec:alps-architektur}

Die Erfüllung der beschriebeben Anforderungen erfordert Modifikationen an der bestehenden Architektur von \gls{lps}. Dafür übernimmt \gls{alps} die gleiche vorteilhafte asynchrone Architektur von \gls{lps} mit Modifikationen an den benutzten Technologien sowie Transportprotokollen, um Kontrolle über den Fluss der Alarmmeldungen zu gewährleisten und Koordination zwischen den Komponenten zu ermöglichen. Zusätzlich werden fehlende Funktionalitäten bei \gls{lps} wie \textit{Idempotenz} durch weitere Komponenten ergänzt. Als Resultat ergibt sich eine neue modifizierte Architektur, die aus vier Kernkomponenten besteht:

\begin{itemize}
    \item \textbf{\gls{alpone}}: Übernimmt die Aufgaben von \gls{lpone}. Wird hinter \gls{srv}~1 positioniert. Empfängt Alarmmeldungen von \gls{srv}~1 und leitet sie an den Message Broker weiter.
    \item \textbf{Message Broker}: Speichert die eingehenden Alarmmeldungen zwischen und verteilt sie an \acrshort{alptwo}.
    \item \textbf{\gls{alptwo}}: Übernimmt die Aufgaben von \gls{lptwo}. Wird vor \gls{srv}~2 positioniert. Konsumiert Alarmmeldungen vom Message Broker und leitet sie an \gls{srv}~2 weiter.
    \item \textbf{\gls{alpdedup}}: Implementiert die Deduplizierungslogik von Alarmmeldungen. Wird unmittelbar vor dem letzten Service in der Verarbeitungskette positioniert (\gls{srv}~3 im \gls{marc}-System). Empfängt Alarmmeldungen von vorgelagerten Services und leitet nur die eindeutigen Alarmmeldungen an den nachgelagerten Service weiter.
\end{itemize}

Als Transportprotokoll zwischen den Komponenten wird \gls{amqp} verwendet, da \gls{amqp} die benötigten Mechanismen bietet, die die Kontrolle über den Fluss der Nachrichten bzw.~Alarmmeldungen und die Koordination zwischen den Komponenten ermöglichen.

\subsection{In-Order-Delivery Garantie}
\label{sub-sec:alps-in-order-delivery-garantie}

Um die die \textit{In-Order-Delivery} zu garantieren, sollen mithilfe von \gls{amqp} Mechanismen folgende Anforderungen erfüllt werden:

\begin{itemize}
    \item Alarmmeldungen pro Alarmanlage sollen immer von einer bestimmten \gls{lptwo} Instanz verarbeitet werden: Beliebige Verteilung von Alarmmeldungen auf beliebige \gls{lptwo} Instanzen soll unterbunden werden.
    \item Alarmmeldungen pro Alarmanlage sollen immer sequentiell von einer einzigen \gls{lptwo} Instanz verarbeitet werden: Parallele Verarbeitung von Alarmmeldungen einer einzelnen Alarmanlage soll unterbunden werden. Jedoch ist die parallele Verarbeitung von Alarmmeldungen verschiedener Alarmanlagen weiterhin erlaubt.
    \item Übertragungsfehler in \gls{lptwo} Instanzen sollen nicht zu einer Verlust führen: Alarmmeldungen, die von einer \gls{lptwo} Instanz nicht verarbeitet werden können, sollen neu übertragen werden, bis sie erfolgreich verarbeitet wurden. Dabei soll die Reihenfolge der Alarmmeldungen pro Alarmanlage beibehalten werden.
\end{itemize}

Die Erfüllung dieser Anforderungen garantiert die \textit{In-Order-Delivery} pro Alarmanlage.

\subsection{Exactly-Once-Delivery Garantie}
\label{sub-sec:alps-exactly-once-delivery-garantie}